\scanpage{4}
\begin{center}\uline{Formules de Néron}\end{center}

$A_K$ V.A. sur $K$ corps des fractions de $V$, anneau
de val. discrète (corps résiduel fini \unsure{\}|?} ?)

On définit une fonction $\unsure{v|\nu}(X,\underline a)$ \inserted{à valeurs dans $\mathbb Q$} sur la partie
de $D_a(\deleted{\unsure{A}}A_K)\times\unsure{\mathfrak Z|Z}_0^*(A_K)\deleted{\unsure{K}}$ formée des couples $X,\underline a$ tels
que $\operatorname{supp}X\cap\operatorname{supp}\underline a=\varnothing$, uniquement déterminée
par les propriétés suivantes :

(i) bilinéaire.

(ii) Invariant par translation par les $a\in A(K)$

(iii) Si $X\underset{\ell}{\sim}0$, \deleted{Alors} \deleted{et} $X=\operatorname{div}(f)$, on a

\[
 v(X,\underline a)=\sum_{x_0\in A_0}i(\deleted{\unsure{\overline{\phantom{x}}}}\operatorname{div}_{\unsure{s_0|f_0|g}}(\bar f),\overline{\underline a}\,;x_0)
\]

% slanted
\margin{Si $m=1$ i.e. $A_0$ connexe, \uline{alors} $v(X,\underline a)=0$ si $X\underset{\rm alg}{\sim}0$ et par \unsure{suite|vraiment} $v(X,\underline a)_V=0$ pour \unsure{tt|$\forall$} $X,\underline a$}

$A$ est le \uline{modèle de Néron} de $A$, $A_0
=A\otimes_V k$, $\bar f$ est $f$ considérée comme fonction sur $A$,
\deleted{et} $\operatorname{div}_{s_0}\bar f$ est la composante verticale de $\operatorname{div}(\bar f)$,
$\overline{\underline a}$ est l’``adhérence'' de $\underline a$ dans $A$. \}

(iv) Pour $X\in D_a(A_K)$ fixé, $0\notin\operatorname{supp}X$, et
$x\in A(K)$, $x\notin\operatorname{supp}X$ variable, $v(X,(x)-(0))$ est
borné.

De plus, on a les propriétés suivantes :

(iv bis) $v(X,(x)-(a))$ ne dépend (pour $X\in D_a(A_K)$
\deleted{fixé, on} et $a\in A(K)$, $a\notin\operatorname{supp}X$ fixés) \deleted{$x$ variable}
que de la classe de $x$ mod $A(K)^0={}$Im. inv. de $A_0^0(\unsure{k})$
donc ne prend qu’un nb fini de valeurs.
% archive sheet number: 3
